%% ====================================================================
%% Rapport au format article, deux colonnes
%% Méthode de recherche en informatique (MRI), FASI, UPC
%% Travail pratique n° 2 : comparaison de performance entre technologies
%%
%% Compilation : pdfLaTeX (Overleaf : laisser le compilateur par défaut)
%%   pdflatex rapport.tex   (deux fois, pour les renvois)
%% Les fichiers figures/pgf_*.tex et tables/*.tex sont générés par
%% analysis/analyse.py (valeurs lues dans raw-results/) puis copiés par
%% final-report/sync.py : ne pas les modifier à la main.
%% ====================================================================
\documentclass[conference,a4paper]{IEEEtran}

\usepackage[T1]{fontenc}
\usepackage[utf8]{inputenc}
\usepackage[french]{babel}
\usepackage{mathptmx}
\usepackage{graphicx}
\usepackage{booktabs}
\usepackage{amsmath}
\usepackage{url}
\usepackage{pgfplots}
\pgfplotsset{compat=1.18}
% Ajouts au modèle : boîtes à moustaches, compatibilité TikZ/babel (« ; »
% actif en français) et nombres affichés à la française (virgule décimale).
\usepgfplotslibrary{statistics}
\usetikzlibrary{babel}
\pgfkeys{/pgf/number format/.cd, use comma, 1000 sep={\,}}
\usepackage[hidelinks]{hyperref}

% Le modèle IEEE place la légende des tableaux en petites capitales :
% on conserve ce comportement par défaut.

\renewcommand{\IEEEkeywordsname}{Mots-clés}
\renewcommand{\abstractname}{Résumé}

\begin{document}

%% --------------------------------------------------------------------
%% TITRE
%% --------------------------------------------------------------------
\title{Laravel, FastAPI, Spring Boot ou Gin~: que vaut vraiment la
performance d'une API REST quand la base de données sature~?}

%% --------------------------------------------------------------------
%% AUTEURS
%% --------------------------------------------------------------------
\author{
\IEEEauthorblockN{Victoire TSHIMANGA MBIA}
\IEEEauthorblockA{FASI, UPC\\
{[adresse@email.com]}}
\and
\IEEEauthorblockN{Benel ILOKO LOKULI}
\IEEEauthorblockA{FASI, UPC\\
{[adresse@email.com]}}
\and
\IEEEauthorblockN{Jureny AMANI}
\IEEEauthorblockA{FASI, UPC\\
{[adresse@email.com]}}
\and
\IEEEauthorblockN{Moise MABIALA MODIMBI}
\IEEEauthorblockA{FASI, UPC\\
{[adresse@email.com]}}
\and
\IEEEauthorblockN{[Prénom NOM Postnom]}
\IEEEauthorblockA{FASI, UPC\\
{[adresse@email.com]}}
}

\maketitle

%% --------------------------------------------------------------------
%% RÉSUMÉ
%% --------------------------------------------------------------------
\begin{abstract}
Le choix d'un framework web est souvent justifié par des classements de
performance obtenus sur des cas très simples, loin d'une application
réelle qui interroge une base de données. L'objectif de cette étude est
de mesurer si la technologie change réellement le débit, la latence et
la consommation de ressources d'une même API REST. Nous comparons
Laravel (PHP), FastAPI (Python), Spring Boot (Java) et Gin (Go), qui
implémentent à l'identique l'API d'un mini réseau social adossée à
MySQL~8.4. Chaque version, limitée à deux cœurs et 1~Go de mémoire, a
subi un test de charge k6 à 10, 50 et 100~utilisateurs simultanés,
répété 30~fois (360~essais), en mesurant le débit, les latences p50,
p95 et p99, le taux d'erreur, le CPU, la mémoire et le temps de
démarrage. MySQL sature dès 10~utilisateurs pour les quatre versions~:
le débit reste compris entre 384 et 485~requêtes/s et l'écart entre
technologies ne dépasse pas 20~\%, sans aucune erreur. Les différences
se trouvent ailleurs~: Gin sert ce débit avec environ 22~\% d'un cœur
et 11~Mo, Laravel sature ses deux cœurs, Spring Boot occupe environ
210~Mo et met 3,3~s à démarrer, contre 0,45~s pour Gin. Pour une
application limitée par sa base de données, la technologie pèse donc
surtout sur le coût d'exécution, et peu sur le débit.
\end{abstract}

\begin{IEEEkeywords}
performance des API REST, test de charge, frameworks web, MySQL
\end{IEEEkeywords}

%% ====================================================================
\section{Introduction}
%% ====================================================================
La performance d'une application web compte pour ses utilisateurs, qui
attendent une réponse rapide, et pour ceux qui l'exploitent, qui paient
le processeur et la mémoire consommés. Elle devient critique dans les
situations de forte charge~: pics d'audience, montée en charge
automatique, hébergement sur des machines modestes.

Des affirmations circulent sur la performance des langages et des
frameworks~: \og Go est plus rapide que PHP \fg, \og Java consomme
trop de mémoire \fg. Elles reposent surtout sur des classements publics
comme TechEmpower~\cite{techempower}, qui mesurent des cas
volontairement simples (sérialisation JSON, une seule requête SQL) avec
des implémentations très optimisées. Ces classements ne disent pas ce
qui se passe pour une application ordinaire dont le travail principal
est d'interroger une base de données.

Notre question de recherche est la suivante~: pour une même API REST
adossée à MySQL et écrite de façon équivalente, la technologie
utilisée change-t-elle significativement le débit, la latence et la
consommation de ressources sous une charge croissante~?

Nous avons formulé trois hypothèses avant la campagne de mesures.
H1~: Gin et Spring Boot ont un débit supérieur à FastAPI et Laravel à
100~utilisateurs, car le code compilé (Go) ou compilé à la volée (JVM)
s'exécute plus vite que le code interprété. H2~: sur les écritures, les
écarts se réduisent, car MySQL devient le facteur limitant. H3~: Spring
Boot consomme le plus de mémoire et met le plus de temps à démarrer, en
raison du coût de la machine virtuelle Java.

Cet article apporte quatre implémentations vérifiées comme
équivalentes, un protocole de mesure reproductible et 360~essais dont
les données brutes sont publiées. La section~\ref{sec:methodo} décrit la
méthodologie, la section~\ref{sec:dev} le développement des versions,
la section~\ref{sec:resultats} les résultats et leur discussion, avant
la conclusion.

%% ====================================================================
\section{Méthodologie}\label{sec:methodo}
%% ====================================================================
\subsection{Application de référence}
L'application est l'API d'un mini réseau social. Elle offre cinq
fonctionnalités~: lister les publications, paginées, avec leurs nombres
de mentions \og j'aime \fg{} et de commentaires (\texttt{GET
/api/posts})~; créer une publication~; aimer une publication~; suivre
un utilisateur~; commenter une publication. Les données manipulées sont
réparties en cinq tables MySQL (utilisateurs, publications, mentions
\og j'aime \fg, abonnements, commentaires). Le jeu initial compte
1\,000~utilisateurs, 10\,000~publications, 20\,000~mentions
\og j'aime \fg, 5\,000~abonnements et 20\,000~commentaires, générés avec
une graine aléatoire fixe. Le comportement attendu est fixé par un
contrat commun~: requête SQL exacte de chaque route, ordre des
vérifications, codes HTTP (200, 201, 400, 404, 409) et corps JSON.
Par exemple, une mention \og j'aime \fg{} en double renvoie le code 409
avec le corps \texttt{\{"error":"already\_liked"\}}.

\subsection{Technologies comparées}
Les quatre technologies du tableau~\ref{tab:techno} ont été choisies
parce qu'elles représentent quatre modèles d'exécution courants~: des
processus PHP qui reconstruisent leur contexte à chaque requête
(Laravel avec php-fpm et nginx), une boucle d'événements asynchrone en
Python (FastAPI avec uvicorn), une machine virtuelle avec compilation à
la volée et un pool de threads (Spring Boot avec Tomcat) et un binaire
natif avec des goroutines (Gin). La base est MySQL~8.4.13 (InnoDB).

\begin{table}[htbp]
\caption{Technologies comparées}
\label{tab:techno}
\centering
\begin{tabular}{@{}lll@{}}
\toprule
\textbf{Technologie} & \textbf{Langage} & \textbf{Version} \\
\midrule
Laravel & PHP 8.3.36 & 13.42.0 \\
FastAPI & Python 3.12.16 & 0.163.1 \\
Spring Boot & Java (OpenJDK 25.0.5) & 4.1.2 \\
Gin & Go 1.27.4 & 1.12.0 \\
\bottomrule
\end{tabular}
\end{table}

\subsection{Métriques}
Nous jugeons qu'une version est plus performante qu'une autre à l'aide
des grandeurs suivantes.
\begin{itemize}
\item \textbf{Débit}~: nombre de requêtes traitées par seconde
(req/s), mesuré par k6~; il exprime la capacité de traitement
soutenue.
\item \textbf{Latence p50, p95 et p99}~: temps de réponse en
dessous duquel se trouvent 50, 95 et 99~\% des requêtes (ms), mesuré
par k6~; la médiane décrit l'expérience typique, p95 et p99 les
ralentissements subis par une partie des utilisateurs.
\item \textbf{Taux d'erreur}~: part des requêtes en échec (\%),
mesurée par k6~; un bon débit obtenu avec des erreurs n'est pas un bon
résultat. Un code 409 (doublon) est une réponse normale et ne compte
pas comme une erreur.
\item \textbf{CPU moyen de l'API}~: en \% d'un cœur (200~\% = deux
cœurs), relevé environ une fois par seconde par \texttt{docker stats}~;
il mesure le coût d'exécution. Pour Laravel, php-fpm et nginx sont
additionnés.
\item \textbf{RAM maximale de l'API}~: en Mo, relevée par
\texttt{docker stats}~; elle mesure l'empreinte mémoire.
\item \textbf{CPU de MySQL}~: en \%, relevé par \texttt{docker stats}~;
il indique si le facteur limitant est la base ou l'API.
\item \textbf{Efficacité}~: débit divisé par le CPU consommé
(req/s par cœur), calculée~; elle compare les versions à ressources
égales.
\item \textbf{Temps de démarrage à froid}~: en ms, du lancement du
conteneur à la première réponse 200~; il compte pour la montée en
charge automatique.
\end{itemize}

\subsection{Outils et environnement de test}
La charge est générée par k6 2.4.0~\cite{k6}. Le CPU et la mémoire sont
relevés par \texttt{docker stats}, et l'ensemble est piloté par des
scripts Python. Une seule machine sert à la fois de serveur et de
générateur de charge~: un ordinateur portable HP Victus (Intel Core
i5-12450H, 8~cœurs et 12~threads, 16~Go de RAM) sous Windows~11, avec
Docker Desktop 29.9.3 (WSL~2). Chaque API tourne dans un conteneur
limité à 2~cœurs et 1~Go~; le conteneur nginx de Laravel dispose en
plus de 0,5~cœur et 128~Mo. MySQL est limité à 2~cœurs et 2~Go, avec un
cache InnoDB de 512~Mo. Chaque API ouvre au plus 10~connexions MySQL et
accède à la base en SQL direct, sans ORM ni cache. La durabilité de
MySQL est réduite (\texttt{innodb\_flush\_log\_at\_trx\_commit=2}, sans
journal binaire~\cite{mysql}). Avec la configuration par défaut, chaque
écriture coûtait environ 9~ms de synchronisation disque sous WSL~2, et
les mesures auraient porté sur le disque plutôt que sur les API.

\subsection{Déroulement des tests}
Les tests suivent l'ordre ci-dessous, entièrement automatisé par
\texttt{bench/run\_campaign.py}.
\begin{enumerate}
\item Chaque essai utilise le même scénario k6 à nombre constant
d'utilisateurs virtuels (VU). Chaque VU enchaîne les requêtes sans
pause et tire chacune au sort~: liste 50~\%, création de publication
15~\%, mention \og j'aime \fg{} 15~\%, abonnement 10~\%,
commentaire~10~\%. Les identifiants sont tirés dans le jeu initial et
les pages entre 1 et~50, de 20~publications.
\item La campagne compte trois tours de dix essais. L'ordre des
technologies est décalé d'un cran à chaque tour.
\item Pour chaque technologie, seuls ses conteneurs sont démarrés (les
autres API sont arrêtées). Le script attend une réponse 200, puis
échauffe l'API pendant 10~s à 10~VU.
\item Pour chaque essai et chaque charge (10, 50 et 100~VU), la base
est remise dans son état initial et les comptages sont vérifiés. Après
une pause de 2~s, k6 tourne pendant 15~s, et \texttt{docker stats} est
relevé en parallèle.
\item Le démarrage à froid est mesuré à part
(\texttt{bench/startup\_time.py}), 30~fois par API~: de
\texttt{docker compose up -d} à la première réponse 200 sur
\texttt{GET /api/posts?limit=1} (sondage toutes les 50~ms), suivi d'un
relevé de la mémoire après 5~s de repos.
\end{enumerate}
On obtient 4~technologies $\times$ 3~charges $\times$ 30~essais
$=$ 360~essais, exécutés en 1~h~46 le 15~novembre~2026 sans
interruption. Pour chaque configuration, nous calculons la médiane, la
moyenne et un intervalle de confiance à 95~\% sur les 30~essais, comme
le recommandent Georges et al.~\cite{georges}. Les différences de p95
entre technologies sont testées par Kruskal-Wallis~\cite{kruskal},
sans hypothèse de normalité, puis par paires par
Mann-Whitney~\cite{mannwhitney} avec la correction de Holm~\cite{holm}.

%% ====================================================================
\section{Développement}\label{sec:dev}
%% ====================================================================
\subsection{Réalisation des versions}
Chaque version tient en un fichier principal de 250 à 420~lignes~: un
contrôleur pour Laravel et pour Spring Boot, \texttt{main.py} pour
FastAPI et \texttt{main.go} pour Gin. Les bibliothèques d'accès aux
données sont \texttt{DB::select} et \texttt{DB::insert} (PDO) pour
Laravel, \texttt{aiomysql} pour FastAPI, \texttt{JdbcTemplate} avec le
pool HikariCP pour Spring Boot et \texttt{database/sql} avec le pilote
\texttt{go-sql-driver/mysql} pour Gin.

Trois points méritent d'être signalés. D'abord, la validation des
entrées est écrite à la main dans les quatre versions, car les
validateurs fournis par les frameworks n'appliquent pas les mêmes
règles~: un entier envoyé sous forme de chaîne (\texttt{"12"}) est
accepté par certains et refusé par d'autres. Ensuite, les dates sont
formatées par MySQL (\texttt{DATE\_FORMAT}) et renvoyées telles
quelles. Enfin, chaque version est configurée pour la production~:
mode debug coupé, journaux d'accès désactivés, OPcache et caches de
configuration pour Laravel, \texttt{ReleaseMode} pour Gin,
\texttt{-Xmx512m} pour Spring Boot. Les versions des dépendances sont
figées.

L'assistant Claude Code (Anthropic) a été utilisé pour écrire le code
des quatre versions, les scripts de mesure et d'analyse, ainsi qu'une
première version de ce rapport. Le groupe a fixé le cahier des charges,
tranché chaque point ambigu et validé chaque étape~; toutes les valeurs
citées proviennent des fichiers de mesures bruts.

\subsection{Équivalence fonctionnelle}
Nous avons vérifié l'équivalence par un test de contrat unique
(\texttt{tests/contract\_test.py}) de 102~cas, exécuté sur les quatre
API. Il couvre les cas nominaux, les erreurs de validation (champ
manquant, mauvais type, entier en chaîne, texte vide ou trop long), les
ressources absentes, les doublons, l'auto-abonnement, les bornes de la
pagination, le tri et le traitement des espaces Unicode. Pour chaque
cas, il contrôle le code HTTP, le corps de la réponse et l'état de la
base. Il compare ensuite le corps exact des réponses des quatre API,
hors identifiants et dates des créations. Les quatre API réussissent
les 102~cas, avec des corps identiques dans tous les cas. Ce résultat
était la condition préalable à toute mesure.

\subsection{Difficultés rencontrées}
Les formats d'erreur par défaut différaient~: code 422 pour Laravel et
FastAPI, page \og Whitelabel \fg{} pour Spring Boot, redirection sur la
barre oblique finale pour Gin. Nous les avons remplacés par des
gestionnaires qui renvoient tous \texttt{\{"error":"\ldots"\}}. Les
pilotes MySQL de Laravel et de Gin préparaient chaque requête côté
serveur (préparation, exécution, fermeture), contrairement à ceux de
FastAPI et de Spring Boot. Les quatre versions insèrent désormais les
paramètres côté client. Le connecteur de Laravel envoie en plus
\texttt{USE} et \texttt{SET NAMES} à chaque requête HTTP~; ce
comportement d'origine est conservé et signalé dans les limites. Le
coût des écritures, d'environ 9~ms avec la durabilité complète, a
conduit au réglage décrit plus haut. Enfin, la structure du résumé JSON
de k6 en version~2 a été vérifiée sur un essai court avant l'écriture
du script d'extraction.

%% ====================================================================
\section{Résultats et discussion}\label{sec:resultats}
%% ====================================================================
\subsection{Données recueillies}
Les 360~essais sont tous présents (30 par technologie et par charge),
sans valeur manquante, et représentent 2\,294\,821~requêtes HTTP. Aucune
requête n'a échoué~: le taux d'erreur est de 0~\% dans tous les essais.
Chaque essai produit, dans le fichier CSV \texttt{mesures.csv}, une
ligne par route et une ligne globale. Il est accompagné du résumé JSON
de k6 et des relevés de \texttt{docker stats}. Les 120~démarrages à
froid sont enregistrés dans \texttt{demarrage.csv}. Pour contrôler la
cohérence des données, nous avons appliqué la loi de
Little~\cite{little}. Avec des utilisateurs qui enchaînent les requêtes
sans pause, elle impose $N = X \cdot R$ (nombre de VU $=$ débit
$\times$ latence moyenne). Le rapport VU/débit et la latence moyenne
mesurée concordent à 1,1~\% près dans les 12~configurations. Le
tableau~\ref{tab:resultats} présente la synthèse des mesures.

\begin{table*}[t]
\caption{Synthèse des mesures (médianes sur 30~essais de 15~s)}
\label{tab:resultats}
\centering
% Généré par analysis/analyse.py -- ne pas modifier à la main.
\begin{table*}[t]
\centering
\small
\caption{Résultats globaux (tous endpoints) : médiane sur 30 essais de 15~s}
\label{tab:global}
\begin{tabular}{llrrrrrrrr}
\toprule
API & VU & req/s & p50 (ms) & p95 (ms) & p99 (ms) & CPU API (\%) & RAM (Mo) & CPU MySQL (\%) & req/s/cœur \\
\midrule
\multirow{3}{*}{Laravel} & 10 & 403,6 & 13,9 & 70,3 & 77,1 & 198 & 58 & 198 & 204 \\
 & 50 & 383,5 & 114,8 & 184,0 & 197,1 & 194 & 59 & 200 & 192 \\
 & 100 & 383,8 & 274,5 & 305,2 & 369,8 & 191 & 60 & 200 & 199 \\
\midrule
\multirow{3}{*}{FastAPI} & 10 & 485,4 & 9,4 & 74,5 & 80,4 & 58 & 107 & 191 & 839 \\
 & 50 & 444,2 & 100,8 & 290,3 & 401,3 & 57 & 107 & 189 & 778 \\
 & 100 & 441,5 & 201,2 & 603,8 & 930,9 & 56 & 107 & 199 & 775 \\
\midrule
\multirow{3}{*}{Spring Boot} & 10 & 448,0 & 8,9 & 81,2 & 86,1 & 39 & 211 & 191 & 1\,156 \\
 & 50 & 404,7 & 101,0 & 292,8 & 398,8 & 40 & 212 & 196 & 1\,002 \\
 & 100 & 406,2 & 199,2 & 601,8 & 900,3 & 40 & 215 & 200 & 1\,023 \\
\midrule
\multirow{3}{*}{Gin} & 10 & 449,2 & 8,4 & 83,2 & 87,5 & 23 & 11 & 200 & 1\,933 \\
 & 50 & 422,6 & 99,9 & 287,1 & 383,0 & 21 & 11 & 192 & 1\,957 \\
 & 100 & 420,5 & 197,6 & 593,3 & 800,3 & 22 & 13 & 192 & 1\,935 \\
\bottomrule
\end{tabular}
\par\smallskip\footnotesize CPU en \% d'un cœur (200~\% = limite de 2 cœurs). Laravel : php-fpm + nginx sommés (limite 2,5 cœurs). RAM : maximum observé pendant l'essai. Efficacité = débit / (CPU API / 100).
\end{table*}

\end{table*}

\subsection{Analyse}
La figure~\ref{fig:mysql} montre que le CPU de MySQL est à
189--200~\%, c'est-à-dire au maximum de ses deux cœurs, pour les quatre
API et dès 10~VU. Au même moment, trois des quatre API utilisent moins
de 60~\% d'un seul cœur (figure~\ref{fig:cpu}). Le facteur limitant est
donc la base de données.

\begin{figure}[htbp]
\centering
% Généré par analysis/analyse.py -- ne pas modifier à la main.
\begin{tikzpicture}
\begin{axis}[
  width=\columnwidth, height=5cm,
  xlabel={Utilisateurs virtuels simultanés (VU)},
  ylabel={CPU moyen de MySQL (\%)},
  xtick={10,50,100},
  tick label style={font=\footnotesize},
  label style={font=\footnotesize},
  legend style={font=\scriptsize, at={(0.97,0.03)}, anchor=south east},
  ymin=0,
  ymax=250]
\addplot[black, thick, mark=*, error bars/.cd, y dir=both, y explicit]
  coordinates {(10,195.53) +- (0,2.07) (50,196.97) +- (0,1.90) (100,198.98) +- (0,1.51)};
\addplot[black, thick, dashed, mark=square*, error bars/.cd, y dir=both, y explicit]
  coordinates {(10,192.26) +- (0,2.47) (50,191.50) +- (0,1.87) (100,194.76) +- (0,2.26)};
\addplot[black, thick, dotted, mark=triangle*, error bars/.cd, y dir=both, y explicit]
  coordinates {(10,193.66) +- (0,1.94) (50,195.29) +- (0,2.02) (100,196.70) +- (0,1.90)};
\addplot[gray, thick, dashdotted, mark=diamond*, error bars/.cd, y dir=both, y explicit]
  coordinates {(10,196.06) +- (0,1.91) (50,194.69) +- (0,2.08) (100,194.49) +- (0,1.92)};
\legend{Laravel, FastAPI, Spring Boot, Gin}
\end{axis}
\end{tikzpicture}

\caption{CPU moyen de MySQL selon la charge (moyenne et IC~95~\% sur
30~essais). La limite de ses deux cœurs est de 200~\%.}
\label{fig:mysql}
\end{figure}

La figure~\ref{fig:debit} montre que le débit plafonne entre 378 et
485~req/s (moyennes). Il baisse légèrement entre 10 et 50~VU, puis ne varie plus
à 100~VU. FastAPI a le débit le plus élevé à chaque charge (médiane de
441,5~req/s à 100~VU), devant Gin (420,5), Spring Boot (406,2) et
Laravel (383,8)~: l'écart entre le premier et le dernier est de
15~\%.

\begin{figure}[htbp]
\centering
% Généré par analysis/analyse.py -- ne pas modifier à la main.
\begin{tikzpicture}
\begin{axis}[
  width=\columnwidth, height=5cm,
  xlabel={Utilisateurs virtuels simultanés (VU)},
  ylabel={Débit global (req/s)},
  xtick={10,50,100},
  tick label style={font=\footnotesize},
  label style={font=\footnotesize},
  legend style={font=\scriptsize, at={(0.97,0.03)}, anchor=south east},
  ymin=0]
\addplot[black, thick, mark=*, error bars/.cd, y dir=both, y explicit]
  coordinates {(10,401.73) +- (0,4.10) (50,378.43) +- (0,5.60) (100,379.28) +- (0,5.15)};
\addplot[black, thick, dashed, mark=square*, error bars/.cd, y dir=both, y explicit]
  coordinates {(10,484.20) +- (0,6.29) (50,441.36) +- (0,3.97) (100,438.18) +- (0,5.04)};
\addplot[black, thick, dotted, mark=triangle*, error bars/.cd, y dir=both, y explicit]
  coordinates {(10,440.96) +- (0,7.02) (50,401.75) +- (0,3.97) (100,403.09) +- (0,5.31)};
\addplot[gray, thick, dashdotted, mark=diamond*, error bars/.cd, y dir=both, y explicit]
  coordinates {(10,447.46) +- (0,3.34) (50,420.65) +- (0,2.46) (100,419.06) +- (0,3.26)};
\legend{Laravel, FastAPI, Spring Boot, Gin}
\end{axis}
\end{tikzpicture}

\caption{Débit global selon la charge (moyenne et IC~95~\% sur
30~essais).}
\label{fig:debit}
\end{figure}

La figure~\ref{fig:p95} montre que la latence p95 augmente presque
proportionnellement à la charge pour FastAPI, Spring Boot et Gin~:
d'environ 80~ms à 10~VU à environ 600~ms à 100~VU. Celle de Laravel
monte moins vite et atteint 305~ms à 100~VU.

\begin{figure}[htbp]
\centering
% Généré par analysis/analyse.py -- ne pas modifier à la main.
\begin{tikzpicture}
\begin{axis}[
  width=\columnwidth, height=5cm,
  xlabel={Utilisateurs virtuels simultanés (VU)},
  ylabel={Latence p95 globale (ms)},
  xtick={10,50,100},
  tick label style={font=\footnotesize},
  label style={font=\footnotesize},
  legend style={font=\scriptsize, at={(0.03,0.97)}, anchor=north west},
  ymin=0]
\addplot[black, thick, mark=*, error bars/.cd, y dir=both, y explicit]
  coordinates {(10,70.30) +- (0,0.30) (50,186.77) +- (0,2.26) (100,316.52) +- (0,9.65)};
\addplot[black, thick, dashed, mark=square*, error bars/.cd, y dir=both, y explicit]
  coordinates {(10,74.34) +- (0,0.51) (50,293.79) +- (0,6.64) (100,613.95) +- (0,10.71)};
\addplot[black, thick, dotted, mark=triangle*, error bars/.cd, y dir=both, y explicit]
  coordinates {(10,81.37) +- (0,0.22) (50,293.65) +- (0,0.95) (100,610.58) +- (0,9.72)};
\addplot[gray, thick, dashdotted, mark=diamond*, error bars/.cd, y dir=both, y explicit]
  coordinates {(10,83.21) +- (0,0.09) (50,286.87) +- (0,1.08) (100,594.91) +- (0,3.01)};
\legend{Laravel, FastAPI, Spring Boot, Gin}
\end{axis}
\end{tikzpicture}

\caption{Latence p95 globale selon la charge (moyenne et IC~95~\% sur
30~essais).}
\label{fig:p95}
\end{figure}

À 50~VU, Laravel a la latence médiane la plus élevée (114,8~ms contre
99,9 à 101,0~ms pour les autres), mais le p95 le plus bas (184,0~ms
contre 287,1 à 292,8~ms) et le p99 le plus bas (197,1~ms contre 383,0
à 401,3~ms). La figure~\ref{fig:box} montre que ses 30~essais ne
recouvrent pas ceux des autres API. La figure~\ref{fig:routes} montre
que cet écart se retrouve sur chacune des cinq routes.

\begin{figure}[htbp]
\centering
% Généré par analysis/analyse.py -- ne pas modifier à la main.
\begin{tikzpicture}
\begin{axis}[
  width=\columnwidth, height=5cm,
  boxplot/draw direction=y,
  xtick={1,2,3,4},
  xticklabels={Laravel,FastAPI,Spring Boot,Gin},
  ylabel={Latence p95 globale (ms), 50~VU},
  tick label style={font=\footnotesize},
  label style={font=\footnotesize}]
\addplot+[boxplot prepared={draw position=1, lower whisker=181.07, lower quartile=182.91, median=184.03, upper quartile=189.97, upper whisker=203.27}, draw=black, fill=gray!40, solid] coordinates {};
\addplot+[boxplot prepared={draw position=2, lower whisker=284.17, lower quartile=287.83, median=290.31, upper quartile=292.46, upper whisker=385.33}, draw=black, fill=gray!40, solid] coordinates {};
\addplot+[boxplot prepared={draw position=3, lower whisker=290.89, lower quartile=291.90, median=292.75, upper quartile=294.46, upper whisker=300.26}, draw=black, fill=gray!40, solid] coordinates {};
\addplot+[boxplot prepared={draw position=4, lower whisker=278.11, lower quartile=286.02, median=287.08, upper quartile=288.16, upper whisker=293.64}, draw=black, fill=gray!40, solid] coordinates {};
\end{axis}
\end{tikzpicture}

\caption{Distribution du p95 global des 30~essais à 50~VU (boîte~:
quartiles et médiane~; moustaches~: minimum et maximum).}
\label{fig:box}
\end{figure}

\begin{figure}[htbp]
\centering
% Généré par analysis/analyse.py -- ne pas modifier à la main.
\begin{tikzpicture}
\begin{axis}[
  width=\columnwidth, height=5cm,
  ybar, bar width=4pt, ymin=0,
  symbolic x coords={GET posts,POST posts,like,follow,comment,global},
  xtick=data,
  x tick label style={rotate=30, anchor=east},
  ylabel={Latence p95 (ms), 50~VU},
  tick label style={font=\footnotesize},
  label style={font=\footnotesize},
  legend style={font=\scriptsize, at={(0.5,0.97)}, anchor=north, legend columns=4},
  enlarge x limits=0.1,
  enlarge y limits={upper, value=0.25}]
\addplot[fill=gray!80, draw=black, error bars/.cd, y dir=both, y explicit]
  coordinates {(GET posts,189.67) +- (0,2.37) (POST posts,182.42) +- (0,2.20) (like,181.93) +- (0,2.14) (follow,182.13) +- (0,2.38) (comment,182.89) +- (0,2.28) (global,186.77) +- (0,2.26)};
\addplot[fill=gray!50, draw=black, error bars/.cd, y dir=both, y explicit]
  coordinates {(GET posts,296.60) +- (0,6.85) (POST posts,287.47) +- (0,7.66) (like,290.67) +- (0,7.04) (follow,291.02) +- (0,6.28) (comment,289.52) +- (0,5.64) (global,293.79) +- (0,6.64)};
\addplot[fill=gray!20, draw=black, error bars/.cd, y dir=both, y explicit]
  coordinates {(GET posts,287.69) +- (0,1.63) (POST posts,295.62) +- (0,1.43) (like,296.81) +- (0,0.92) (follow,294.51) +- (0,1.65) (comment,303.25) +- (0,2.86) (global,293.65) +- (0,0.95)};
\addplot[fill=white, draw=black, error bars/.cd, y dir=both, y explicit]
  coordinates {(GET posts,192.45) +- (0,0.58) (POST posts,294.16) +- (0,0.88) (like,295.18) +- (0,0.90) (follow,294.93) +- (0,1.09) (comment,331.64) +- (0,11.46) (global,286.87) +- (0,1.08)};
\legend{Laravel, FastAPI, Spring Boot, Gin}
\end{axis}
\end{tikzpicture}

\caption{Latence p95 par route à 50~VU (moyenne et IC~95~\% sur
30~essais).}
\label{fig:routes}
\end{figure}

Le tableau~\ref{tab:kruskal} confirme que le p95 diffère selon la
technologie pour chaque route ($p < 10^{-6}$), avec une taille d'effet
$\varepsilon^2$ comprise entre 0,64 et 0,83. Seule la paire
Spring Boot / Gin n'est pas significativement différente, sur la
création de publication et l'abonnement. Avec 30~essais très réguliers,
des écarts de quelques millisecondes deviennent
\og significatifs \fg~: c'est la taille de l'écart qui compte, plus que
la valeur de~$p$.

\begin{table}[htbp]
\caption{Test de Kruskal-Wallis sur le p95 à 50~VU (Mann-Whitney et
correction de Holm pour les paires, $\alpha = 0{,}05$)}
\label{tab:kruskal}
\centering
% Généré par analysis/analyse.py -- ne pas modifier à la main.
\begin{tabular}{@{}lcccp{2.9cm}@{}}
\toprule
\textbf{Route} & \textbf{$H$} & \textbf{$p$} & \textbf{$\varepsilon^2$} & \textbf{Paires non signif.} \\
\midrule
GET posts & 96,8 & $<10^{-6}$ & 0,81 & aucune \\
POST posts & 88,8 & $<10^{-6}$ & 0,74 & Spring Boot / Gin \\
like & 85,4 & $<10^{-6}$ & 0,71 & aucune \\
follow & 77,8 & $<10^{-6}$ & 0,64 & Spring Boot / Gin \\
comment & 99,5 & $<10^{-6}$ & 0,83 & aucune \\
global & 92,1 & $<10^{-6}$ & 0,77 & aucune \\
\bottomrule
\end{tabular}

\end{table}

La figure~\ref{fig:cpu} montre que, pour un débit voisin, le coût en
processeur varie fortement~: à 50~VU, Gin utilise en moyenne 22~\%
d'un cœur, Spring Boot 42~\%, FastAPI 57~\%, tandis que Laravel
sature ses deux cœurs (196~\%, nginx compris). Rapportée au CPU,
l'efficacité de Gin atteint environ 1\,950~req/s par cœur, contre
environ 1\,000 pour Spring Boot, 780 pour FastAPI et 200 pour Laravel
(tableau~\ref{tab:resultats}). La figure~\ref{fig:ram} montre que
Spring Boot occupe le plus de mémoire (212~Mo en moyenne à 50~VU),
devant FastAPI (110~Mo, deux processus), Laravel (64~Mo, dix processus
php-fpm et nginx) et Gin (11~Mo).

\begin{figure}[htbp]
\centering
% Généré par analysis/analyse.py -- ne pas modifier à la main.
\begin{tikzpicture}
\begin{axis}[
  width=\columnwidth, height=5cm,
  ybar, bar width=12pt, ymin=0,
  symbolic x coords={Laravel,FastAPI,Spring Boot,Gin},
  xtick=data,
  ylabel={CPU moyen de l'API (\%), 50~VU},
  nodes near coords={\pgfmathprintnumber[fixed,precision=0]{\pgfplotspointmeta}},
  nodes near coords style={font=\scriptsize},
  tick label style={font=\footnotesize},
  label style={font=\footnotesize},
  enlarge x limits=0.18,
  enlarge y limits={upper, value=0.15}]
\addplot[fill=gray!40, draw=black] coordinates
  {(Laravel,196.19) (FastAPI,57.29) (Spring Boot,42.50) (Gin,21.77)};
\end{axis}
\end{tikzpicture}

\caption{CPU moyen de l'API à 50~VU (moyenne sur 30~essais~; Laravel~:
php-fpm et nginx additionnés).}
\label{fig:cpu}
\end{figure}

\begin{figure}[htbp]
\centering
% Généré par analysis/analyse.py -- ne pas modifier à la main.
\begin{tikzpicture}
\begin{axis}[
  width=\columnwidth, height=5cm,
  ybar, bar width=12pt, ymin=0,
  symbolic x coords={Laravel,FastAPI,Spring Boot,Gin},
  xtick=data,
  ylabel={RAM maximale de l'API (Mo), 50~VU},
  nodes near coords={\pgfmathprintnumber[fixed,precision=0]{\pgfplotspointmeta}},
  nodes near coords style={font=\scriptsize},
  tick label style={font=\footnotesize},
  label style={font=\footnotesize},
  enlarge x limits=0.18,
  enlarge y limits={upper, value=0.15}]
\addplot[fill=gray!40, draw=black] coordinates
  {(Laravel,63.71) (FastAPI,110.11) (Spring Boot,211.87) (Gin,11.28)};
\end{axis}
\end{tikzpicture}

\caption{RAM maximale de l'API pendant un essai à 50~VU (moyenne sur
30~essais).}
\label{fig:ram}
\end{figure}

Enfin, la figure~\ref{fig:demarrage} et le tableau~\ref{tab:demarrage}
montrent que Gin démarre en 445~ms (médiane), Laravel en 902~ms,
FastAPI en 1\,733~ms et Spring Boot en 3\,346~ms. Laravel est le plus
irrégulier (de 806 à 1\,887~ms). Au repos, Gin occupe 3,3~Mo et Spring
Boot 123,2~Mo. L'image de Gin (32~Mo) est 25~fois plus petite que celle
de Laravel avec nginx (808~Mo).

\begin{figure}[htbp]
\centering
% Généré par analysis/analyse.py -- ne pas modifier à la main.
\begin{tikzpicture}
\begin{axis}[
  width=\columnwidth, height=5cm,
  ybar, bar width=12pt, ymin=0,
  symbolic x coords={Laravel,FastAPI,Spring Boot,Gin},
  xtick=data,
  ylabel={Démarrage à froid (ms)},
  nodes near coords={\pgfmathprintnumber[fixed,precision=0]{\pgfplotspointmeta}},
  nodes near coords style={font=\scriptsize},
  tick label style={font=\footnotesize},
  label style={font=\footnotesize},
  enlarge x limits=0.18,
  enlarge y limits={upper, value=0.15}]
\addplot[fill=gray!40, draw=black] coordinates
  {(Laravel,902.25) (FastAPI,1733.40) (Spring Boot,3345.70) (Gin,444.70)};
\end{axis}
\end{tikzpicture}

\caption{Temps de démarrage à froid (médiane sur 30~démarrages).}
\label{fig:demarrage}
\end{figure}

\begin{table}[htbp]
\caption{Démarrage à froid, RAM au repos et taille des images}
\label{tab:demarrage}
\centering
% Généré par analysis/analyse.py -- ne pas modifier à la main.
\begin{tabular}{@{}lcccc@{}}
\toprule
\textbf{Technologie} & \textbf{Médiane} & \textbf{Moyenne} & \textbf{RAM repos} & \textbf{Image} \\
 & ms & ms $\pm$ IC 95~\% & Mo & Mo \\
\midrule
Laravel & 902 & 1\,064 $\pm$ 122 & 46,6 & 808 \\
FastAPI & 1\,733 & 1\,764 $\pm$ 74 & 104,2 & 225 \\
Spring Boot & 3\,346 & 3\,523 $\pm$ 147 & 123,2 & 497 \\
Gin & 445 & 472 $\pm$ 22 & 3,3 & 32 \\
\bottomrule
\end{tabular}

\end{table}

\subsection{Interprétation et discussion}
Les écarts de débit s'expliquent d'abord par la saturation de MySQL~:
le débit mesure surtout la capacité de la base, et les écarts restants,
de 20~\% au plus, ne reflètent pas la vitesse propre des langages. Comme
le débit ne varie plus au-delà de 10~VU, chaque utilisateur
supplémentaire allonge la file d'attente devant la base, ce qui
explique la croissance presque proportionnelle de la latence (loi de
Little). Nous n'avons pas mesuré quelle requête consomme le plus de
CPU dans MySQL. La liste paginée, qui compte les mentions et les
commentaires de chaque publication et saute jusqu'à 980~lignes, est
une candidate plausible, mais non vérifiée.

Pour la queue de latence plus courte de Laravel, nous proposons une
explication que nous n'avons pas vérifiée. php-fpm dispose
d'exactement dix processus, chacun gardant sa connexion MySQL pendant
toute la requête~: au plus dix requêtes sont traitées à la fois, et les
autres attendent dans la file du serveur, servies dans l'ordre
d'arrivée. Chaque requête attend donc à peu près autant que les
autres, d'où une médiane plus haute mais une queue courte. Les trois
autres API acceptent toutes les requêtes en même temps et les font
attendre à l'intérieur du processus, sans ordre garanti, d'où une
dispersion plus grande. Limiter la concurrence dans ces trois versions
permettrait de tester cette explication.

Les résultats confirment en partie nos hypothèses. H1 est infirmée~:
à 100~VU, FastAPI a le meilleur débit, et Gin et Spring Boot ne
dépassent que Laravel. H2 ne peut pas être testée telle qu'elle est
formulée~: le scénario mélange lectures et écritures dans un même
essai, si bien que la base saturée ralentit toutes les routes, et les
écarts entre technologies sont du même ordre pour les lectures et pour
les écritures (figure~\ref{fig:routes}). H3 est confirmée~: Spring Boot
occupe environ deux fois plus de mémoire que FastAPI et 17 à 19~fois
plus que Gin, et c'est la version la plus lente à démarrer.

Ces données ne permettent pas de classer les langages en général~:
elles comparent quatre implémentations dans une situation où la base
limite tout. Avec une base plus puissante, une requête plus légère ou
des réglages poussés (Laravel Octane, réglage fin de la JVM), le
classement du débit pourrait changer. Les écarts de coût en CPU et en
mémoire ont peu de chances de s'inverser, mais nous ne l'avons pas
vérifié.

Notre étude a plusieurs limites. (1)~k6, les API et MySQL partagent
une seule machine, un ordinateur portable, sous Docker Desktop et
WSL~2. (2)~La durabilité de MySQL est réduite~: une panne de l'hôte
pourrait faire perdre environ une seconde de transactions, sans effet
sur les mesures. (3)~Laravel dispose de 0,5~cœur de plus pour nginx et
exécute deux commandes SQL supplémentaires par requête HTTP.
(4)~Chaque essai dure 15~s. (5)~Le scénario, sans temps de pause,
n'imite pas des utilisateurs réels, et un seul mélange de requêtes a
été testé. (6)~Le CPU et la mémoire sont relevés environ une fois par
seconde. (7)~Les configurations sont raisonnables pour la production,
mais pas optimisées au maximum. (8)~Le temps de démarrage inclut le
coût de la commande \texttt{docker compose}~; le conteneur est relancé
et non recréé, et le cache disque de l'hôte reste chaud.

%% ====================================================================
\section{Conclusion}
%% ====================================================================
Nous nous sommes demandé si, pour une même API REST adossée à MySQL, la
technologie change significativement le débit, la latence et la
consommation de ressources. Sur 360~essais, la réponse est nuancée~: le
débit varie peu (20~\% au plus), car MySQL sature dès 10~utilisateurs
simultanés. En revanche, le coût d'exécution varie beaucoup, d'environ
dix fois en CPU entre Gin et Laravel et de 17 à 19~fois en mémoire
entre Gin et Spring Boot, tout comme la forme de la distribution des
latences et le temps de démarrage. Ces conclusions restent limitées à
une seule machine, à un seul scénario et à des configurations
raisonnables mais non optimisées. Pour aller plus loin, il faudrait
refaire la campagne avec une base dimensionnée pour ne plus saturer,
tester séparément lectures et écritures, et vérifier l'explication
proposée pour la queue de latence de Laravel.

%% ====================================================================
%% RÉFÉRENCES
%% ====================================================================
\begin{thebibliography}{00}
\bibitem{techempower} TechEmpower, « Web Framework Benchmarks ». [En
ligne]. Disponible~: \url{https://www.techempower.com/benchmarks/}.
[Consulté le 16~novembre~2026].
\bibitem{k6} Grafana Labs, « Grafana k6 documentation ». [En ligne].
Disponible~: \url{https://grafana.com/docs/k6/latest/}. [Consulté le
16~novembre~2026].
\bibitem{mysql} Oracle, « MySQL 8.4 Reference Manual -- InnoDB Startup
Options and System Variables ». [En ligne]. Disponible~:
\url{https://dev.mysql.com/doc/refman/8.4/en/innodb-parameters.html}.
[Consulté le 16~novembre~2026].
\bibitem{georges} A.~Georges, D.~Buytaert et L.~Eeckhout,
« Statistically rigorous Java performance evaluation »,
\textit{Proc. ACM OOPSLA}, p.~57-76, 2007.
\bibitem{kruskal} W.~H.~Kruskal et W.~A.~Wallis, « Use of ranks in
one-criterion variance analysis », \textit{Journal of the American
Statistical Association}, vol.~47, n\textsuperscript{o}~260,
p.~583-621, 1952.
\bibitem{mannwhitney} H.~B.~Mann et D.~R.~Whitney, « On a test of
whether one of two random variables is stochastically larger than the
other », \textit{The Annals of Mathematical Statistics}, vol.~18,
n\textsuperscript{o}~1, p.~50-60, 1947.
\bibitem{holm} S.~Holm, « A simple sequentially rejective multiple test
procedure », \textit{Scandinavian Journal of Statistics}, vol.~6,
n\textsuperscript{o}~2, p.~65-70, 1979.
\bibitem{little} J.~D.~C.~Little, « A proof for the queuing formula:
$L = \lambda W$ », \textit{Operations Research}, vol.~9,
n\textsuperscript{o}~3, p.~383-387, 1961.
\end{thebibliography}

\end{document}
